\documentclass[a4paper]{article}

%% Language and font encodings
\usepackage[spanish]{babel}
\usepackage[utf8]{inputenc}
\usepackage[T1]{fontenc}
\usepackage{lmodern}
\usepackage{tabularx}
\usepackage[none]{hyphenat}
\usepackage[table]{xcolor}

%% Sets page size and margins
\usepackage[a4paper,top=3cm,bottom=2cm,left=3cm,right=3cm,marginparwidth=2cm]{geometry}

%% Useful packages
\usepackage{amsmath}
\usepackage{graphicx}
\usepackage[colorlinks=true, allcolors=blue]{hyperref}
\usepackage{amsfonts}

\newtheorem{definicion}{Definición}[section]
\newtheorem{proposicion}{Proposición}[section]

\title{Nuestro maravilloso trabajo para la asignatura de\\
Procesamiento de Imágenes Digitales}
\author{Adrián Gutiérrez Roldán \and José Luis Rabadán Morejón \and Iván Ruiz Guerra}

\begin{document}
\maketitle

\begin{abstract}
 \noindent El resumen debe tener como máximo 250 palabras y ser un único párrafo. El resumen debe describir lo más fielmente posible el trabajo realizado (no el propuesto inicialmente, que podría ser distinto). 
 En esta plantilla se presentan las indicaciones para realizar el trabajo de PID, al mismo tiempo que la estructura que debe tener la documentación a presentar. 

\hspace{1cm}

\noindent \textbf{Palabras clave:} 
%(hasta cinco palabras que clasifiquen el trabajo) 
PID, instrucciones, trabajo en grupo, procesamiento de imágenes.
\end{abstract}


\section{Introducción}

...

\section{Planteamiento teórico}

\subsection{Notación y planteamiento del problema}
\[\mathbb{R}^{H_W \times W_W \times n} = \{\text{matrices W con dimensiones} (H_W, W_W), n\in \mathbb{N} \text{ número de canales} \}\]

Sean $I \in \mathbb{R}^{H_I \times W_I \times 3}$ , $S \in \mathbb{R}^{H_S \times W_S \times 3}$ las matrices que representan las dos imágenes que conforman nuestro espacio de trabajo. Nuestro objetivo final es conseguir una imagen cuyo formato estructural corresponda al de la imagen I, pero con el estilo de la imagen S.

\begin{definicion}
Definimos una deformación $\mathbf{d}$ como una función continua de dos variables sobre el dominio $\Omega \subset \mathbb{R}^2$:

\[
\begin{aligned}
\mathbf{d} : \Omega \subset \mathbb{R}^2 &\longrightarrow \mathbb{R}^2 \\
(x, y) &\longmapsto \mathbf{d}(x, y) = \left( x', \, y' \right)
\end{aligned}
\]

donde $x'$ e $y'$ representan las componentes de la deformación (o desplazamiento) en los ejes $x$ e $y$, respectivamente.
\end{definicion}

\begin{definicion}
Llamaremos \textbf{campo de deformación} $\mathbf{D}$ al campo vectorial de desplazamientos asociado a $\phi$, definido por:
\[
\begin{aligned}
\mathbf{D} : \Omega &\longrightarrow \mathbb{R}^2 \\
(x, y) &\longmapsto \mathbf{D}(x, y) = d(x, y) - (x, y) = \left( D_x(x, y), \, D_y(x, y) \right)
\end{aligned}
\]

donde  $D_x(x, y)$, $D_y(x, y)$ representan los desplazamientos horizontales y verticales ,respectivamente, con respecto a las coordenadas de origen (x,y).
\end{definicion}

Para llevar a cabo la labor de deformación de la imagen de estilo S, se introduce una red neuronal parametrizada en función de $\theta$. $N_\theta(I,S)$ recibe como entrada ambas imágenes involucradas en el proceso y devuelve el campo de deformación asociado a la imagen S para tomar las formas de I.

\[
\mathbf{D} = N_\theta(I,S)
\]

En el caso práctico, $\mathbf{D}$ vendrá dado por una matriz $M_D \in \mathbb{R}^{H_W \times W_W \times 2}$ donde el canal 1 almacenará los $D_x(x,y)$ y el canal 2 contendrá los $D_y(x,y)$

\begin{definicion}
Sea X una variable aleatoria discreta, que toma valores en $D_k = \{x_1,x_2,...,x_k\}$ con distribución de probabilidad $P_X$ se define la entropía de X como:
\[
H_X = E(-\ln(P_X(x)) = \sum_{x_k\in D_k} -\ln(P_X(x_k))P_X(x_k)
\]
\end{definicion}

\noindent La entropia mide el grado de incertidumbre sobre la variable X.

\begin{definicion}
Sean X e Y dos variables aleatorias definimos la información mutua entre ellas como:

\[
IM(X,Y) = H_X - H_{X|Y} 
\]

\noindent Esta magnitud cuantifica cualquier tipo de dependencia entre las variables X e Y. Puede verse como cuán determinada queda una variable en función de la incertidumbre con respecto a la otra.
\end{definicion}

\begin{proposicion}
Sean X e Y dos variables aleatorias. Entonces:
\[
IM(X,Y) = 0 \iff \text{X e Y son independientes} 
\]

\end{proposicion}

Los algoritmos usuales de aprendizaje automático se encargan de minimizar una función de pérdida. Nuestro objetivo es maximizar la información mutua entre la imagen estructural I y la imagen de estilo transformada, luego bastaría con añadir un menos delante de la función a maximizar para que se convierta en una minimización:

\[
L_{IM} = - IM(I, S_{deformada})
\]

\section{Implementación}

...

\section{Experimentación}

...

\section{Manual de usuario}

...

\section{Uso de IA en este trabajo}

...

\section{Conclusiones}

...

\section{Autoevaluación}

...

\section{Tabla de tiempos}

\begin{center}
\begin{tabularx}{\textwidth}{|c|c|c|>{\centering\arraybackslash}X|}
\hline
\rowcolor{gray!20}
Fecha de la actividad
& Tiempo (h)
& Miembro
& Actividad realizada
\\\hline
07/10/2026 & 2 & Iván & Búsqueda de tema\\\hline
07/10/2026 & 2 & Adrián & Búsqueda de tema\\\hline
08/10/2026 & 1,5 & Luis & Búsqueda de tema\\\hline
29/10/2026 & 2,5 & Iván & Lectura comprensiva del artículo\\\hline
29/10/2026 & 1,5 & Iván & Análisis del código del repositorio\\\hline
03/11/2026 & 2 & Luis & Lectura comprensiva del artículo\\\hline
03/11/2026 & 1 & Luis & Análisis del código del repositorio\\\hline
04/11/2026 & 2 & Adrián & Lectura comprensiva del artículo\\\hline
04/11/2026 & 1 & Adrián & Análisis del código del repositorio\\\hline
04/11/2026 & 1 & Iván & Seguimiento 1: Decisión y redacción de la planificación y objetivos\\\hline
04/11/2026 & 1 & Luis & Seguimiento 1: Decisión y redacción de la planificación y objetivos\\\hline
04/11/2026 & 1 & Adrián & Seguimiento 1: Decisión y redacción de la planificación y objetivos\\\hline
07/11/2026 & & Adrián & Memoria: Estructura general de las definiciones matemáticas\\\hline
10/11/2026 & 0,5 & Iván & Notebook: Configuración del entorno y carga de imágenes\\\hline
& & & \\\hline
& & & \\\hline
& & & \\\hline
& & & \\\hline
& & & \\\hline
\end{tabularx}
\end{center}


\begin{thebibliography}{10}


\bibitem{clave:libro}
U. N. Experto, \emph{Un libro que escribí},
Editorial, 1996.

\bibitem{clave:revista}
Y. O. Mismo,
``Título del artículo'',
\emph{Revista Publicación Periódica}, Vol. 17, pp. 1-100, 1997.
\end{thebibliography}

\end{document}